% =============================================================================
% cn-gongwen · base.tex — 各变体共用导言（GB/T 9704—2012 版式参数与字体配置）
%
% 用法：每个变体文件在 \documentclass[linespread=1.81, twoside]{ctexart} 之后
%       % =============================================================================
% cn-gongwen · base.tex — 各变体共用导言（GB/T 9704—2012 版式参数与字体配置）
%
% 用法：每个变体文件在 \documentclass[linespread=1.81, twoside]{ctexart} 之后
%       % =============================================================================
% cn-gongwen · base.tex — 各变体共用导言（GB/T 9704—2012 版式参数与字体配置）
%
% 用法：每个变体文件在 \documentclass[linespread=1.81, twoside]{ctexart} 之后
%       % =============================================================================
% cn-gongwen · base.tex — 各变体共用导言（GB/T 9704—2012 版式参数与字体配置）
%
% 用法：每个变体文件在 \documentclass[linespread=1.81, twoside]{ctexart} 之后
%       \input{base}，再按各自格式编排正文。参数说明见 main.tex 头注与 README。
% =============================================================================
% --- 5.2.1 页边与版心 --------------------------------------------------------
% twoside 下 left=订口(内侧)28mm、right=切口(外侧)26mm，双页自动镜像。
% footskip 使页码一字线上缘距版心下缘约 7mm（7.5）。
\usepackage[a4paper, top=37mm, bottom=35mm, left=28mm, right=26mm,
            footskip=10.5mm]{geometry}
\setlength{\parindent}{2\ccwd}   % 正文每自然段左空二字（7.3.3）
\raggedbottom                    % 禁止竖向胶水拉伸，保持 22 行网格与版头定位稳定

% --- 公文红（6.2：红色油墨达到色谱 Y80%、M80%）------------------------------
\usepackage{xcolor}
\definecolor{gwred}{cmyk}{0, 0.8, 0.8, 0}

% --- 字体：正版公文字体自动检测，缺失则回退 ---------------------------------
% 小标宋（发文机关标志、标题）：无免费授权字体，回退为宋体加粗。
\IfFontExistsTF{FZXiaoBiaoSong-B05S}{%
  \setCJKfamilyfont{gwxbs}{FZXiaoBiaoSong-B05S}%
  \newcommand{\xbsong}{\CJKfamily{gwxbs}}%
}{%
  \IfFontExistsTF{方正小标宋简体}{%
    \setCJKfamilyfont{gwxbs}{方正小标宋简体}%
    \newcommand{\xbsong}{\CJKfamily{gwxbs}}%
  }{%
    \newcommand{\xbsong}{\songti\bfseries}%
  }%
}
% 仿宋_GB2312 / 楷体_GB2312（存在则替换 ctex 默认仿宋/楷体）
\IfFontExistsTF{FangSong_GB2312}{%
  \setCJKfamilyfont{gwfs}{FangSong_GB2312}%
  \renewcommand{\fangsong}{\CJKfamily{gwfs}}%
}{}
\IfFontExistsTF{KaiTi_GB2312}{%
  \setCJKfamilyfont{gwkt}{KaiTi_GB2312}%
  \renewcommand{\kaishu}{\CJKfamily{gwkt}}%
}{}
\renewcommand{\familydefault}{\rmdefault}
\AtBeginDocument{\zihao{3}\fangsong}   % 5.2.2 全文默认 3 号仿宋（5.2.3 行距随 linespread）
\xeCJKsetup{CJKecglue={}}              % 公文习惯：汉字与数字/字母间不加空隙（见式样）
\xeCJKDeclareCharClass{CJK}{"2605}     % ★（密级用，7.2.2）按汉字取自中文字体，拉丁字体无此字形

% --- 7.3.3 结构层次：一、黑体 /（一）楷体；1. 与（1）仿宋随正文排 -----------
% 标题行不加额外行距，保持 22 行网格。
\ctexset{
  section = {
    name      = {,、},
    number    = \chinese{section},
    format    = \heiti\zihao{3},
    aftername = {},
    indent    = 2\ccwd,
    beforeskip = 0pt,
    afterskip  = 0pt,
    afterindent = true,
  },
  subsection = {
    name      = {（,）},
    number    = \chinese{subsection},
    format    = \kaishu\zihao{3},
    aftername = {},
    indent    = 2\ccwd,
    beforeskip = 0pt,
    afterskip  = 0pt,
    afterindent = true,
  },
}

% --- 7.5 页码：4号半角宋体数字、两侧一字线；单页右空一字、双页左空一字 -------
\usepackage{fancyhdr}
\pagestyle{fancy}
\fancyhf{}
% 一字 = 4号汉字宽 14bp（页脚上下文中 \ccwd 不可靠，用字面量）
\newcommand{\gwpagenum}{{\zihao{4}\songti —\hspace{4.2bp}\thepage\hspace{4.2bp}—}}
\fancyfoot[RO]{\gwpagenum\hspace*{14bp}}
\fancyfoot[LE]{\hspace*{14bp}\gwpagenum}
\renewcommand{\headrulewidth}{0pt}

% --- 可调参数（标准未规定的量）----------------------------------------------
\newlength{\gwlogodrop}\setlength{\gwlogodrop}{35mm}  % 机关标志上缘距版心上缘(7.2.4)
\newlength{\gwredrule}\setlength{\gwredrule}{0.5mm}   % 版头红色分隔线粗细(标准未定)
\newcommand{\gwlogosize}{\zihao{-0}}                  % 机关标志字号(7.2.4"醒目美观")

\usepackage{graphicx}                % 印章图片可用（见署名处注释）
\usepackage[hidelinks]{hyperref}   % pdftitle 由各变体文件自行 \hypersetup 设置
，再按各自格式编排正文。参数说明见 main.tex 头注与 README。
% =============================================================================
% --- 5.2.1 页边与版心 --------------------------------------------------------
% twoside 下 left=订口(内侧)28mm、right=切口(外侧)26mm，双页自动镜像。
% footskip 使页码一字线上缘距版心下缘约 7mm（7.5）。
\usepackage[a4paper, top=37mm, bottom=35mm, left=28mm, right=26mm,
            footskip=10.5mm]{geometry}
\setlength{\parindent}{2\ccwd}   % 正文每自然段左空二字（7.3.3）
\raggedbottom                    % 禁止竖向胶水拉伸，保持 22 行网格与版头定位稳定

% --- 公文红（6.2：红色油墨达到色谱 Y80%、M80%）------------------------------
\usepackage{xcolor}
\definecolor{gwred}{cmyk}{0, 0.8, 0.8, 0}

% --- 字体：正版公文字体自动检测，缺失则回退 ---------------------------------
% 小标宋（发文机关标志、标题）：无免费授权字体，回退为宋体加粗。
\IfFontExistsTF{FZXiaoBiaoSong-B05S}{%
  \setCJKfamilyfont{gwxbs}{FZXiaoBiaoSong-B05S}%
  \newcommand{\xbsong}{\CJKfamily{gwxbs}}%
}{%
  \IfFontExistsTF{方正小标宋简体}{%
    \setCJKfamilyfont{gwxbs}{方正小标宋简体}%
    \newcommand{\xbsong}{\CJKfamily{gwxbs}}%
  }{%
    \newcommand{\xbsong}{\songti\bfseries}%
  }%
}
% 仿宋_GB2312 / 楷体_GB2312（存在则替换 ctex 默认仿宋/楷体）
\IfFontExistsTF{FangSong_GB2312}{%
  \setCJKfamilyfont{gwfs}{FangSong_GB2312}%
  \renewcommand{\fangsong}{\CJKfamily{gwfs}}%
}{}
\IfFontExistsTF{KaiTi_GB2312}{%
  \setCJKfamilyfont{gwkt}{KaiTi_GB2312}%
  \renewcommand{\kaishu}{\CJKfamily{gwkt}}%
}{}
\renewcommand{\familydefault}{\rmdefault}
\AtBeginDocument{\zihao{3}\fangsong}   % 5.2.2 全文默认 3 号仿宋（5.2.3 行距随 linespread）
\xeCJKsetup{CJKecglue={}}              % 公文习惯：汉字与数字/字母间不加空隙（见式样）
\xeCJKDeclareCharClass{CJK}{"2605}     % ★（密级用，7.2.2）按汉字取自中文字体，拉丁字体无此字形

% --- 7.3.3 结构层次：一、黑体 /（一）楷体；1. 与（1）仿宋随正文排 -----------
% 标题行不加额外行距，保持 22 行网格。
\ctexset{
  section = {
    name      = {,、},
    number    = \chinese{section},
    format    = \heiti\zihao{3},
    aftername = {},
    indent    = 2\ccwd,
    beforeskip = 0pt,
    afterskip  = 0pt,
    afterindent = true,
  },
  subsection = {
    name      = {（,）},
    number    = \chinese{subsection},
    format    = \kaishu\zihao{3},
    aftername = {},
    indent    = 2\ccwd,
    beforeskip = 0pt,
    afterskip  = 0pt,
    afterindent = true,
  },
}

% --- 7.5 页码：4号半角宋体数字、两侧一字线；单页右空一字、双页左空一字 -------
\usepackage{fancyhdr}
\pagestyle{fancy}
\fancyhf{}
% 一字 = 4号汉字宽 14bp（页脚上下文中 \ccwd 不可靠，用字面量）
\newcommand{\gwpagenum}{{\zihao{4}\songti —\hspace{4.2bp}\thepage\hspace{4.2bp}—}}
\fancyfoot[RO]{\gwpagenum\hspace*{14bp}}
\fancyfoot[LE]{\hspace*{14bp}\gwpagenum}
\renewcommand{\headrulewidth}{0pt}

% --- 可调参数（标准未规定的量）----------------------------------------------
\newlength{\gwlogodrop}\setlength{\gwlogodrop}{35mm}  % 机关标志上缘距版心上缘(7.2.4)
\newlength{\gwredrule}\setlength{\gwredrule}{0.5mm}   % 版头红色分隔线粗细(标准未定)
\newcommand{\gwlogosize}{\zihao{-0}}                  % 机关标志字号(7.2.4"醒目美观")

\usepackage{graphicx}                % 印章图片可用（见署名处注释）
\usepackage[hidelinks]{hyperref}   % pdftitle 由各变体文件自行 \hypersetup 设置
，再按各自格式编排正文。参数说明见 main.tex 头注与 README。
% =============================================================================
% --- 5.2.1 页边与版心 --------------------------------------------------------
% twoside 下 left=订口(内侧)28mm、right=切口(外侧)26mm，双页自动镜像。
% footskip 使页码一字线上缘距版心下缘约 7mm（7.5）。
\usepackage[a4paper, top=37mm, bottom=35mm, left=28mm, right=26mm,
            footskip=10.5mm]{geometry}
\setlength{\parindent}{2\ccwd}   % 正文每自然段左空二字（7.3.3）
\raggedbottom                    % 禁止竖向胶水拉伸，保持 22 行网格与版头定位稳定

% --- 公文红（6.2：红色油墨达到色谱 Y80%、M80%）------------------------------
\usepackage{xcolor}
\definecolor{gwred}{cmyk}{0, 0.8, 0.8, 0}

% --- 字体：正版公文字体自动检测，缺失则回退 ---------------------------------
% 小标宋（发文机关标志、标题）：无免费授权字体，回退为宋体加粗。
\IfFontExistsTF{FZXiaoBiaoSong-B05S}{%
  \setCJKfamilyfont{gwxbs}{FZXiaoBiaoSong-B05S}%
  \newcommand{\xbsong}{\CJKfamily{gwxbs}}%
}{%
  \IfFontExistsTF{方正小标宋简体}{%
    \setCJKfamilyfont{gwxbs}{方正小标宋简体}%
    \newcommand{\xbsong}{\CJKfamily{gwxbs}}%
  }{%
    \newcommand{\xbsong}{\songti\bfseries}%
  }%
}
% 仿宋_GB2312 / 楷体_GB2312（存在则替换 ctex 默认仿宋/楷体）
\IfFontExistsTF{FangSong_GB2312}{%
  \setCJKfamilyfont{gwfs}{FangSong_GB2312}%
  \renewcommand{\fangsong}{\CJKfamily{gwfs}}%
}{}
\IfFontExistsTF{KaiTi_GB2312}{%
  \setCJKfamilyfont{gwkt}{KaiTi_GB2312}%
  \renewcommand{\kaishu}{\CJKfamily{gwkt}}%
}{}
\renewcommand{\familydefault}{\rmdefault}
\AtBeginDocument{\zihao{3}\fangsong}   % 5.2.2 全文默认 3 号仿宋（5.2.3 行距随 linespread）
\xeCJKsetup{CJKecglue={}}              % 公文习惯：汉字与数字/字母间不加空隙（见式样）
\xeCJKDeclareCharClass{CJK}{"2605}     % ★（密级用，7.2.2）按汉字取自中文字体，拉丁字体无此字形

% --- 7.3.3 结构层次：一、黑体 /（一）楷体；1. 与（1）仿宋随正文排 -----------
% 标题行不加额外行距，保持 22 行网格。
\ctexset{
  section = {
    name      = {,、},
    number    = \chinese{section},
    format    = \heiti\zihao{3},
    aftername = {},
    indent    = 2\ccwd,
    beforeskip = 0pt,
    afterskip  = 0pt,
    afterindent = true,
  },
  subsection = {
    name      = {（,）},
    number    = \chinese{subsection},
    format    = \kaishu\zihao{3},
    aftername = {},
    indent    = 2\ccwd,
    beforeskip = 0pt,
    afterskip  = 0pt,
    afterindent = true,
  },
}

% --- 7.5 页码：4号半角宋体数字、两侧一字线；单页右空一字、双页左空一字 -------
\usepackage{fancyhdr}
\pagestyle{fancy}
\fancyhf{}
% 一字 = 4号汉字宽 14bp（页脚上下文中 \ccwd 不可靠，用字面量）
\newcommand{\gwpagenum}{{\zihao{4}\songti —\hspace{4.2bp}\thepage\hspace{4.2bp}—}}
\fancyfoot[RO]{\gwpagenum\hspace*{14bp}}
\fancyfoot[LE]{\hspace*{14bp}\gwpagenum}
\renewcommand{\headrulewidth}{0pt}

% --- 可调参数（标准未规定的量）----------------------------------------------
\newlength{\gwlogodrop}\setlength{\gwlogodrop}{35mm}  % 机关标志上缘距版心上缘(7.2.4)
\newlength{\gwredrule}\setlength{\gwredrule}{0.5mm}   % 版头红色分隔线粗细(标准未定)
\newcommand{\gwlogosize}{\zihao{-0}}                  % 机关标志字号(7.2.4"醒目美观")

\usepackage{graphicx}                % 印章图片可用（见署名处注释）
\usepackage[hidelinks]{hyperref}   % pdftitle 由各变体文件自行 \hypersetup 设置
，再按各自格式编排正文。参数说明见 main.tex 头注与 README。
% =============================================================================
% --- 5.2.1 页边与版心 --------------------------------------------------------
% twoside 下 left=订口(内侧)28mm、right=切口(外侧)26mm，双页自动镜像。
% footskip 使页码一字线上缘距版心下缘约 7mm（7.5）。
\usepackage[a4paper, top=37mm, bottom=35mm, left=28mm, right=26mm,
            footskip=10.5mm]{geometry}
\setlength{\parindent}{2\ccwd}   % 正文每自然段左空二字（7.3.3）
\raggedbottom                    % 禁止竖向胶水拉伸，保持 22 行网格与版头定位稳定

% --- 公文红（6.2：红色油墨达到色谱 Y80%、M80%）------------------------------
\usepackage{xcolor}
\definecolor{gwred}{cmyk}{0, 0.8, 0.8, 0}

% --- 字体：正版公文字体自动检测，缺失则回退 ---------------------------------
% 小标宋（发文机关标志、标题）：无免费授权字体，回退为宋体加粗。
\IfFontExistsTF{FZXiaoBiaoSong-B05S}{%
  \setCJKfamilyfont{gwxbs}{FZXiaoBiaoSong-B05S}%
  \newcommand{\xbsong}{\CJKfamily{gwxbs}}%
}{%
  \IfFontExistsTF{方正小标宋简体}{%
    \setCJKfamilyfont{gwxbs}{方正小标宋简体}%
    \newcommand{\xbsong}{\CJKfamily{gwxbs}}%
  }{%
    \newcommand{\xbsong}{\songti\bfseries}%
  }%
}
% 仿宋_GB2312 / 楷体_GB2312（存在则替换 ctex 默认仿宋/楷体）
\IfFontExistsTF{FangSong_GB2312}{%
  \setCJKfamilyfont{gwfs}{FangSong_GB2312}%
  \renewcommand{\fangsong}{\CJKfamily{gwfs}}%
}{}
\IfFontExistsTF{KaiTi_GB2312}{%
  \setCJKfamilyfont{gwkt}{KaiTi_GB2312}%
  \renewcommand{\kaishu}{\CJKfamily{gwkt}}%
}{}
\renewcommand{\familydefault}{\rmdefault}
\AtBeginDocument{\zihao{3}\fangsong}   % 5.2.2 全文默认 3 号仿宋（5.2.3 行距随 linespread）
\xeCJKsetup{CJKecglue={}}              % 公文习惯：汉字与数字/字母间不加空隙（见式样）
\xeCJKDeclareCharClass{CJK}{"2605}     % ★（密级用，7.2.2）按汉字取自中文字体，拉丁字体无此字形

% --- 7.3.3 结构层次：一、黑体 /（一）楷体；1. 与（1）仿宋随正文排 -----------
% 标题行不加额外行距，保持 22 行网格。
\ctexset{
  section = {
    name      = {,、},
    number    = \chinese{section},
    format    = \heiti\zihao{3},
    aftername = {},
    indent    = 2\ccwd,
    beforeskip = 0pt,
    afterskip  = 0pt,
    afterindent = true,
  },
  subsection = {
    name      = {（,）},
    number    = \chinese{subsection},
    format    = \kaishu\zihao{3},
    aftername = {},
    indent    = 2\ccwd,
    beforeskip = 0pt,
    afterskip  = 0pt,
    afterindent = true,
  },
}

% --- 7.5 页码：4号半角宋体数字、两侧一字线；单页右空一字、双页左空一字 -------
\usepackage{fancyhdr}
\pagestyle{fancy}
\fancyhf{}
% 一字 = 4号汉字宽 14bp（页脚上下文中 \ccwd 不可靠，用字面量）
\newcommand{\gwpagenum}{{\zihao{4}\songti —\hspace{4.2bp}\thepage\hspace{4.2bp}—}}
\fancyfoot[RO]{\gwpagenum\hspace*{14bp}}
\fancyfoot[LE]{\hspace*{14bp}\gwpagenum}
\renewcommand{\headrulewidth}{0pt}

% --- 可调参数（标准未规定的量）----------------------------------------------
\newlength{\gwlogodrop}\setlength{\gwlogodrop}{35mm}  % 机关标志上缘距版心上缘(7.2.4)
\newlength{\gwredrule}\setlength{\gwredrule}{0.5mm}   % 版头红色分隔线粗细(标准未定)
\newcommand{\gwlogosize}{\zihao{-0}}                  % 机关标志字号(7.2.4"醒目美观")

\usepackage{graphicx}                % 印章图片可用（见署名处注释）
\usepackage[hidelinks]{hyperref}   % pdftitle 由各变体文件自行 \hypersetup 设置
